\documentclass[11pt]{article}

\usepackage[french]{babel}
\usepackage[utf8]{inputenc}
\usepackage{fancyhdr}
\usepackage{fourier}
\usepackage{lastpage}
\usepackage{placeins}
\usepackage{subfigure}
\usepackage[pdftex]{graphicx}
\usepackage{float}
% des commandes utiles pour ecrire des maths
\newcommand{\dx}{\,dx}
\newcommand{\dt}{\,dt}
\newcommand{\ito}{,\dotsc,}
\newcommand{\R}{\mathds{R}}
\newcommand{\N}{\mathds{N}}
\newcommand{\Poly}[1]{\mathcal{P}_{#1}}
\newcommand{\abs}[1]{\left\lvert#1\right\rvert}
\newcommand{\norm}[1]{\left\lVert#1\right\rVert}
\newcommand{\pars}[1]{\left(#1\right)}
\newcommand{\bigpars}[1]{\bigl(#1\bigr)}
\newcommand{\set}[1]{\left\{#1\right\}}
\newenvironment{Reponse}[1]{
\par
\setlength\fboxrule{2pt}%
\begin{minipage}{.95\linewidth}
\fbox{\parbox[c]{.972\linewidth}{\textbf{Propriété} \\ \medskip  #1}} 
\end{minipage}
}


\newenvironment{Reponse2}[1]{%
    %\vspace{3pt}
    \par
  \begin{minipage}{.95\linewidth}
    %\vspace{12pt}\par
    \xdef\@formanswerline{\@questionheader}%
    \xdef\@maskanswerline{\@questionheader}%
    \fbox{\parbox[c]{.972\linewidth}{\textbf{Q\thequestion} : #1}}
    \vspace\questionspace\par
    %{\@answerstitlefont\@answerstitle}
    \vspace{-2ex}
    \begin{multicols}{5}
    \begin{list}{\@answernumberfont\Alph{@choice}.~}
    	{\addtolength{\leftmargin}{-3pt}%
    	%\setlength{\labelsep}{-1pt}
    	\usecounter{@choice}}}{%
    \end{list}
    \end{multicols}
    \@initorcheck%
    \addtocontents{frm}{\@formanswerline\protect\\\protect\hline}%
    \addtocontents{msk}{\@maskanswerline\protect\\\protect\hline}%
  \end{minipage}
  }
  

%%%%%%
% Pour mise-en-forme des fichiers Ada
%
% voir exemple en fin de ce fichier.
%
% ATTENTION, requiert encoding utf-8 (voir 2ième "\lstset" ci-dessous)
 
\usepackage{listings}
\lstset{
  morekeywords={abort,abs,accept,access,all,and,array,at,begin,body,
      case,constant,declare,delay,delta,digits,do,else,elsif,end,entry,
      exception,exit,for,function,generic,goto,if,in,is,limited,loop,
      mod,new,not,null,of,or,others,out,package,pragma,private,
      procedure,raise,range,record,rem,renames,return,reverse,select,
      separate,subtype,task,terminate,then,type,use,when,while,with,
      xor,abstract,aliased,protected,requeue,tagged,until,public,static,void},
  sensitive=f,
  morecomment=[l]\#,
  morestring=[d]",
  showstringspaces=false,
  basicstyle=\small\ttfamily,
  keywordstyle=\bf\small,
  commentstyle=\itshape,
  stringstyle=\sf,
  extendedchars=true,
  columns=[c]fixed,
  numbers = left,
  numberstyle=\tiny,
  stepnumber=2,
  numbersep=5pt,
  framexleftmargin=5mm,
  language=JAVA,
  frame=shadowbox
}

% CI-DESSOUS: conversion des caractères accentués UTF-8 
% en caractères TeX dans les listings...
\lstset{
  literate=%
  {À}{{\`A}}1 {Â}{{\^A}}1 {Ç}{{\c{C}}}1%
  {à}{{\`a}}1 {â}{{\^a}}1 {ç}{{\c{c}}}1%
  {É}{{\'E}}1 {È}{{\`E}}1 {Ê}{{\^E}}1 {Ë}{{\"E}}1% 
  {é}{{\'e}}1 {è}{{\`e}}1 {ê}{{\^e}}1 {ë}{{\"e}}1%
  {Ï}{{\"I}}1 {Î}{{\^I}}1 {Ô}{{\^O}}1%
  {ï}{{\"i}}1 {î}{{\^i}}1 {ô}{{\^o}}1%
  {Ù}{{\`U}}1 {Û}{{\^U}}1 {Ü}{{\"U}}1%
  {ù}{{\`u}}1 {û}{{\^u}}1 {ü}{{\"u}}1%
}

%%%%%%%%%%
% TAILLE DES PAGES (A4 serré)

\setlength{\parindent}{0pt}
\setlength{\parskip}{1ex}
\setlength{\textwidth}{17cm}
\setlength{\textheight}{23cm}
\setlength{\oddsidemargin}{-.7cm}
\setlength{\evensidemargin}{-.7cm}
\setlength{\topmargin}{-.5in}

%%%%%%%%%%
% EN-TÊTES ET PIED DE PAGES

\pagestyle{fancyplain}
\renewcommand{\headrulewidth}{0pt}
\addtolength{\headheight}{1.6pt}
\addtolength{\headheight}{2.6pt}
\lfoot{}
\cfoot{}
\rfoot{\footnotesize\sf page~\thepage/\pageref{LastPage}}
\lhead{\footnotesize\sf Projet GL}
\rhead{\footnotesize\sf ENSIMAG-2A GL10} % numéro d'équipe Teide 

% titre, auteur et date
\title{Documentation technique du compilateur DecaC\\
~\\
\includegraphics[scale=0.75]{logo_ensimag.jpg} 
}

%% penser à indiquer: numéro d'équipe Teide + noms + groupe TD
\author{Julien Khayat, Guillaume Sasdy, Simon Perrat, Charles Roussel, Simon Cathébras}
\date{Le \today}
\begin{document}

Description du projet
Git. Wiki. Open source, logiciel libre. Tout le monde peut apporter sa contribution. Passerelle développé l'année dernière(Definition code de base):  éditer un wiki dans navigateur web -> image miroir sous forme de fichiers et répertoires afin de profiter des                fonctionnalités de Git (explicité). Avantages.
Notre équipe: environnement de test. Pourquoi ? Validation par la communauté du code de l'année dernière afin qu'il soit     intégré à la version "officielle" du logiciel Git.
Autre équipe en parallèle qui fonctionnalités (equipe fonctionnalité).
Methode agile. Interraction communauté. Progression des branches de git.
!!!!!! On chie agile !!!!

Objectifs et périmètre
3 objectifs:
*) Acceptation de notre environnement par la communauté. Critère : notre code sur 'branche next'
*) Couverture de test pour le code de l'année dernière. Critère : nombre de fonctionnalités testées et validées par la       communauté Git.
*) Mise en forme de nos tests et respect des normes de la communauté : nombre de d'insultes ...
+indicateurs de perf pour ces objectif

Périmètre. Concerne les manipulations usuelles sur un wiki: ajout, édition, suppression d'une page et vérification de la     cohérence avec le repository git équivalent. (environnement)
test automatique pour passerelle (tests)
Tester le code de base. Technos utilisées : Perl + Shell.
Perimètre étendu : debug du code de base + aider l'équipe fonctionnalité.

Organistation de l'équipe et Plannification.
	Dates = les sprints.
	Preciser qu'on est en agile.
	Interlocuteur = Charles (equipe fonctionnalité).
	Reunion + reporting (cf notes).
	Organisation Scrum : Tableau des User Stories+Avancement+burndown.
	Communication : Avec M.Moy + Avec Communauté git (feedback).
	
	Prevention de risques et dépendances:
		-Manque d'interet du projet par la communauté. S => envoyer vite et 		etre actif.
		-Envoyer les travaux trop tard. S => bien s'organiser : envoi à 			chaque fin de sprint.
		-L'autre équipe depend de nous. S => choix d'un interlocuteur, 				developper l'environnement rapidement
	
BLA BLA de la fin (pipoooooooooooo)

	
\section{Contexte du projet}
Dans ce projet, 
\end{document}